%% Template for SDP report, adapted from mlp_cw2_template, 2018. 

%% Based on  LaTeX template for ICML 2017 - example_paper.tex at 
%%  https://2017.icml.cc/Conferences/2017/StyleAuthorInstructions

\documentclass{article}
\input{mlp2018_includes}
\usepackage{fancyvrb}

%% You probably do not need to change anything above this comment

%% REPLACE the details in the following commands with your details
\setGroupNumber{7}
\setGroupName{Halo}
\setProductName{Halo Assist}
\setDemoNumber{3}
\setLogoFileName{figs/logo.png}

\begin{document} 

\makeSDPTitle{Demo}


\begin{abstract} 
HaloAssist is a wearable haptic guidance system intended to provide object-search assistance and spatial awareness for partially sighted users through real-time vibrotactile feedback.

In Demo 3, the system was consolidated into a complete end-to-end pipeline. The implementation replaces discrete motor triggering with continuous offset-based haptic mapping, integrates local offline speech transcription via Vosk~\cite{alphacep_vosk}, introduces lightweight LLM-based noun extraction from natural-language commands~\cite{gerganov2023llama}, and applies YOLO-World~\cite{cheng2024yoloworld} for open-vocabulary target detection.
\end{abstract}

\section{Project Management Update} 

\subsection{Project Goals}

The primary objective of this demo was to deliver a reliable end-to-end assistive workflow in which a user issues a natural spoken command and receives continuous directional haptic guidance toward a detected target object.

Table~\ref{tab:progress} summarizes progress against the specific objectives.

\begin{table*}[t]
\centering
\setlength{\tabcolsep}{4pt}
\renewcommand{\arraystretch}{1.3}
\begin{tabular}{|p{0.18\textwidth}|p{0.74\textwidth}|}
\hline
\textbf{Status} & \textbf{Objective} \\
\hline
Achieved & 
\begin{minipage}[t]{\linewidth}
\begin{itemize}
    \item Implementing continuous offset-based haptic feedback for precise directional guidance
    \item Integrating local offline Speech-to-Text (STT) parsing with Vosk and an LLM for robust target noun extraction
    \item Utilizing YOLO-World for dynamic, open-vocabulary object detection based on voice commands
    \item Calculating the object's spatial offset relative to the camera to direct the motor intensity
\end{itemize}
\end{minipage} \\
\hline
Partially Achieved & 
\begin{minipage}[t]{\linewidth}
\begin{itemize}
    \item Miniaturizing hardware components into the final 3D-printed CAD enclosure
    \item End-to-end latency optimization for real-time responsiveness on the Raspberry Pi 5~\cite{rpi5}
    \item Expanding to a full 8-motor haptic array setup
\end{itemize}
\end{minipage} \\
\hline
\end{tabular}
\caption{Progress Overview}
\label{tab:progress}

\vspace{3mm}
\refstepcounter{table}\label{tab:team}
\textbf{Table \thetable.} Team Contributions and Responsibilities
\vspace{2mm}

\begin{tabular}{|p{0.16\textwidth}|p{0.18\textwidth}|p{0.08\textwidth}|p{0.50\textwidth}|}
\hline
\textbf{Team} & \textbf{Member} & \textbf{Hours} & \textbf{Role/Contribution} \\
\hline
Hardware & Member 1 & 40 & Led the hardware team in integrating the I2C MUX + DRV2605 and continued 3D modelling the system enclosure. \\
\hline
Perception & Member 2 & 42 & Integrated YOLO-World dynamic class updating and optimised standard YOLO pipeline for inference on the Pi. \\
\hline
Haptics & Member 3 & 38 & Developed and verified the continuous offset motor mapping, dead zone logic, and motor pulse interval behaviour. \\
\hline
Software & Member 4 & 35 & Implemented the local Vosk STT pipeline, integrated the LLM noun extraction, and managed the main system state logic handling button interrupts. \\
\hline
\end{tabular}
\end{table*}

\subsection{Deviations from Intended Goals}

\textbf{Hardware packaging:} CAD and enclosure iterations progressed substantially; however, full miniaturization remains incomplete due to repeated integration dependencies among camera placement, cable routing, and haptic-driver mounting.

\textbf{Latency tuning:} The perception loop is functionally stable, but end-to-end latency remains under optimization to maintain consistent performance on Raspberry Pi 5 across varying illumination and object classes.

\textbf{NLP pipeline design:} The initial direct-label parsing approach was replaced by a two-stage pipeline in which Vosk transcribes full utterances and a fine-tuned LLM extracts the target noun phrase. This change improved robustness for natural commands such as ``find my water bottle''.

\subsection{Team and System Organization}

HaloAssist is an end-to-end assistive system composed of tightly integrated hardware and software components. The workflow couples a physical trigger interface, speech capture, natural-language target extraction, vision-based localization, and continuous directional haptic output. During this demo cycle, development focused on maturing each subsystem while preserving full-pipeline operability for repeated evaluations.

From earlier demonstrations to the current version, four tracks were advanced in parallel: (i) hardware reliability and enclosure iteration, (ii) perception and model integration, (iii) continuous haptic control logic, and (iv) orchestration tooling for reproducible execution. The resulting architecture supports command-driven object search by extracting a target noun phrase from speech, applying target-specific YOLO-World detection, and mapping visual offset to continuous motor intensity.

In addition, a visualizer website was developed to support both debugging and demonstration (refer to Appendix C, Figure~\ref{fig:website} for the user interface). This interface exposes detections, system states, and runtime behavior during integration tests, thereby reducing diagnosis time for cross-module issues.

To maximize efficiency and collaboration within the team, we used the following working practices:
\begin{itemize}
    \item \textbf{Task tracking and ownership:} We maintained feature ownership by subsystem (Hardware, Perception, Haptics, Software) and used GitHub~\cite{github} issue/task tracking to prioritize integration-critical work. To ensure structural oversight across these tracks, a centralized timeline was established and maintained via Google Docs (see Appendix B, Figure~\ref{fig:timeline}).
    \item \textbf{Regular coordination:} We conducted weekly sync meetings and short checkpoint updates before integration milestones to surface blockers early.
    \item \textbf{Code integration discipline:} We used branch-based development with pull requests and peer review before merging to reduce regression risk in the main loop.
    \item \textbf{Bash-based reproducibility:} We used Bash scripts for repeatable setup and experiment execution across machines (development and Raspberry Pi) to keep runs consistent and reduce manual configuration errors.
    \item \textbf{Verification-driven iteration:} We validated changes with dedicated scripts for camera, detection, speech, latency, and haptic mapping, then iterated parameters based on measured output.
    \item \textbf{Platform profile consistency:} We used profile-based configuration (Pi3, Pi4, Pi5, Mac) so the same codebase can be tested locally and deployed to target hardware with minimal changes.
\end{itemize}

In the next phase, the same organizational model will be retained for final-stage optimization, with emphasis on end-to-end latency reduction, improved ergonomics, and expanded command coverage while preserving guidance stability.

\subsection{Team Workload}

The current workload allocation and ownership structure are summarized in Table~\ref{tab:team}. Values remain provisional and will be finalized in the final submission.

\subsection{Budget Summary}

Current prototype expenditure is primarily driven by compute hardware and haptics electronics. Cost reduction in the next iteration is expected through enclosure consolidation, wiring simplification, and bulk procurement of repeated components.

\section{Quantitative Analysis and Testing}

The evaluation protocol was strictly designed to assess the core functionalities of the HaloAssist pipeline. The primary areas of investigation include (i) the reliability and precision of continuous haptic guidance, (ii) the responsiveness and end-to-end latency of the speech-to-detection loop processing, and (iii) the practical system behavior and stability under realistic operational conditions. 

To ensure the system meets the requirements for real-time assistive applications, quantitative metrics were gathered across varying environmental conditions and interaction patterns. This section details the testing methodologies and presents the empirical results obtained from subsystem validation.

\subsection{Continuous Haptic Mapping Tests}

The continuous offset mapping implementation, utilizing an I2C multiplexer~\cite{nxp_i2c} and DRV2605 haptic motor drivers~\cite{ti_drv2605} in real-time playback mode, was systematically evaluated across the full horizontal camera field of view. The goal was to quantify the smoothness and interpretability of the directional feedback.

\begin{itemize}
    \item \textbf{Relevance:} The accuracy of continuous offset estimation, the stability of pulse-interval behavior, and the strict enforcement of dead-zone boundaries were assessed. These metrics are critical to ensure that the directional feedback provided to the user is both precise and unambiguously interpretable without causing sensory overload.
    \item \textbf{Interpretation:} Experimental data confirmed that the mapping successfully produced a smooth, continuous strength gradient ranging from -1.0 (far left) to +1.0 (far right). To prevent distracting vibrations when the target is centered, a dead zone of $\pm 12\%$ relative to the frame center was implemented, which effectively suppressed center-region noise and flickering. Furthermore, the motor pulse intervals correctly varied from $0.45\text{\,s}$ (slow) to $0.55\text{\,s}$ (fast) in direct correlation with the computed feedback intensity.
    \item \textbf{Robustness:} A rigorous automated sweep over 21 distinct horizontal positions within a 1280px resolution frame was performed to identify any non-linearities or signal discontinuities. Using a predefined jump threshold of $< 0.3$ to flag irregular motor actuation, zero abrupt transitions were observed across all test instances. This indicates a highly stable and predictable linear scaling behavior across the entire visual field.
\end{itemize}

\begin{table}[h]
\centering
\small
\setlength{\tabcolsep}{3pt}
\renewcommand{\arraystretch}{1.2}
\begin{tabular}{p{0.28\columnwidth}p{0.22\columnwidth}p{0.26\columnwidth}p{0.12\columnwidth}}
\hline
\textbf{Test Category} & \textbf{Metric / Objective} & \textbf{Range / Result} & \textbf{Status} \\
\hline
Continuous Offset mapping & Jump threshold $< 0.3$ & Smooth gradient from -1.0 to 1.0 & $\surd$ PASS \\
Pulse Interval Behavior & Fast/Slow mapping & 1.82 to 2.22 pulses/sec & $\surd$ PASS \\
Dead Zone Stability & Prevent center flickering & $\pm$ 12\% (152px width) inactive & $\surd$ PASS \\
\hline
\end{tabular}
\caption{Continuous Motor Mapping Test Results}
\label{tab:haptic-tests}
\end{table}

\subsection{Speech and Detection Pipeline Timing}

This subsystem was comprehensively tested to verify low operational latency. Table~\ref{tab:pipeline-latency} logs the detailed breakdown components of the interaction loop under simulated use conditions.

\begin{figure}[h]
    \centering
    % \includegraphics[width=0.8\linewidth]{figs/latency_plot.png}
    \vspace{2cm} % Placeholder space
    \caption{Latency distribution breakdown showing primary processing bottlenecks vs. system target limits.}
    \label{fig:latency-distribution}
\end{figure}

The latency data indicate that the system architecture meets operational benchmarks, yet isolated processing spikes during LLM NLP extraction and YOLO-World edge deployment remain targets for optimization. These measurements emphasize the ongoing challenges of real-time multi-modal inference on embedded platforms like the Raspberry Pi 5.

\begin{table}[h]
\centering
\small
\setlength{\tabcolsep}{3pt}
\renewcommand{\arraystretch}{1.2}
\begin{tabular}{p{0.36\columnwidth}p{0.30\columnwidth}p{0.22\columnwidth}}
\hline
\textbf{Pipeline Step} & \textbf{Current (ms)} & \textbf{Target (ms)} \\
\hline
Button press to listen start & 120 & 100 \\
Vosk transcription & 380 & 300 \\
LLM noun extraction (llama.cpp) & 95 & 70 \\
YOLO-World first valid detection & 210 & 160 \\
Haptic update command issue & 45 & 30 \\
End-to-end command to guidance & 850 & 650 \\
\hline
\end{tabular}
\caption{Speech-to-guidance latency breakdown (draft values).}
\label{tab:pipeline-latency}
\end{table}

\section{Budget}
Table~\ref{tab:budget-current} provides a draft cost breakdown for the current prototype build. These values are provisional and will be replaced with exact purchase totals in the final version.

\begin{table}[h]
\centering
\small
\setlength{\tabcolsep}{3pt}
\renewcommand{\arraystretch}{1.2}
\begin{tabular}{p{0.42\columnwidth}p{0.20\columnwidth}p{0.22\columnwidth}}
\hline
\textbf{Item} & \textbf{Qty} & \textbf{Cost (\pounds)} \\
\hline
Raspberry Pi 5 + storage + power & 1 & 115 \\
Camera module and mounting & 1 & 35 \\
DRV2605 drivers + I2C multiplexer & 1 set & 28 \\
Haptic motors and wiring materials & 1 set & 32 \\
3D printing materials and enclosure parts & 1 set & 40 \\
Auxiliary electronics and connectors & 1 set & 20 \\
\hline
\textbf{Total (draft)} &  & \textbf{270} \\
\hline
\end{tabular}
\caption{Current prototype budget breakdown (draft values).}
\label{tab:budget-current}
\end{table}

In the next revision, cost will be reduced through compact enclosure redesign and consolidation of repeated electronics purchases.

\section{Miscellaneous}
\label{sec:misc}

This section details several components crucial to the overall efficacy, usability, and architecture of the HaloAssist project. This includes insights drawn from human-computer interaction studies, diagnostic tools built to facilitate development, software architectural decisions, integration of applied Natural Language Processing pipelines, and the underpinning mathematical models controlling real-time spatial haptic interpretation.

\subsection{User Study and Feedback Mechanisms}

Iterative user-facing trials conducted throughout the integration phase formed the core feedback loop for system refinement. The testing protocol was primarily designed to evaluate the intuitive phrasing required for voice commands, the ergonomic comfort and physical weight load of the wearable components, and the cognitive clarity of continuous haptic guidance mechanisms.

Study participants were explicitly instructed to utilize natural, grammatically unstructured voice commands—simulating real-world scenarios rather than strict label matching strategies. While following the generated haptic cues to align their orientation with physical target objects, qualitative feedback and navigation errors were logged.

The most critical findings synthesized from the study were:
\begin{itemize}
    \item Participants exhibited a strong preference for complex, sentence-style commands over rigid keyword interactions, necessitating the adoption of a sophisticated noun-extraction mechanism within the NLP loop.
    \item Qualitative feedback indicated that the continuous modulation of haptic intensity was cognitively easier to interpret and trace than traditional binary left/right discrete vibration bursts. The continuous approach allowed users to estimate radial distance more accurately.
    \item Ergonomic observations underscored that wearability, thermal management of the Raspberry Pi 5 under load, and precise mounting comfort remain primary constraints preventing extended operational sessions. Current prototype weight distribution specifically impacts long-term wearability.
\end{itemize}

These empirical observations heavily influenced the subsequent architectural restructuring and prioritized parameters for system configuration tuning, emphasizing computational offloading where possible to simplify embedded load on the Raspberry Pi 5~\cite{rpi5}.

\subsection{Diagnostic Observability Interface}

To alleviate iterative testing friction, a lightweight, dedicated Web-based observability portal was developed and integrated locally into the pipeline to streamline monitoring processes. This architecture allowed real-time telemetry inspection inside the unified feedback loop, operating as an essential diagnostic layer during demonstrations and field testing.

The functional scope of the visualizer interface encompasses:
\begin{itemize}
    \item Rendering live camera frames overlaid with YOLO-World detection bounding boxes and confidence scores.
    \item Publishing real-time systemic state transitions (e.g., transitions between `listening`, `transcribing`, `extracting targets`, `visual scanning`, and `haptic actuation`). 
    \item Displaying live inference timing data and target spatial confidence signals to aid root-cause diagnosis of latency or localization variations.
\end{itemize}

The deployment of this diagnostic mechanism drastically reduced overall debugging turnaround time and vastly enhanced the coordination capabilities across the perception, haptic, and hardware teams during end-to-end integration stress tests.

\subsection{Systemic Subsystem Architecture}

HaloAssist utilizes a tightly coupled, closed-loop assistive computational architecture to facilitate sequential information processing:
\begin{enumerate}
    \item A mechanical hardware interrupt (button press) asynchronously triggers the command listening phase.
    \item A localized, offline implementation of Vosk STT~\cite{alphacep_vosk} transcribes the raw captured audio into unstructured text strings.
    \item A locally-hosted Large Language Model (LLM), optimized specifically for low-latency edge inference, parses the transcription to extract the highest-probability semantic target noun or noun phrase.
    \item The extracted target is piped as dynamic class label input into the YOLO-World~\cite{cheng2024yoloworld} object detection model via inference node queries.
    \item Based on the visual bounding box coordinates generated, the software calculates an offset value. This spatial target position is transformed mathematically and output directly to the I2C-driven continuous haptic motor array.
\end{enumerate}

By localizing all processing logic and model configurations to the hardware itself, the architecture mitigates latency vulnerabilities previously presented by network dependence, rendering the device fully standalone.

\subsection{NLP Integration using LLM and Llama.cpp}

Historically, matching noisy and unstructured Speech-to-Text outputs natively to localized fixed class labels manifested critical inaccuracies. To correct this, a fine-tuned Large Language Model acting as a semantic parser replaced naive string-matching algorithms. This shift vastly improved system tolerance and accuracy when handling natural human colloquialisms, complex phrasing, embedded filler words, and subjective adjectival modifiers.

Given the rigid compute and memory constraints imposed by the Raspberry Pi 5 platform, achieving viable inference rates required integration of \texttt{llama.cpp}~\cite{gerganov2023llama}. The \texttt{llama.cpp} execution framework optimizes hardware acceleration and mitigates the otherwise unacceptable parsing overhead typical of transformer architectures while fully preserving linguistic target-extraction fidelity.

\subsection{Continuous Haptic Control Mapping}

The implemented mapping in the haptic controller is continuous and position-dependent.

For object center $x_t$, frame center $x_c$, and frame width $W$:
\[
o = \text{clip}\left(\frac{x_t - x_c}{W/2}, -1, 1\right), \quad m = |o|
\]

If the object is near center ($m < b$, where $b$ is the center band), both motors provide a soft cue:
\[
I_L = I_R = I_c
\]

Otherwise, dominant intensity is:
\[
I_d = m^{\gamma} \cdot I_{\max}, \quad I_s = r \cdot I_d
\]
where $\gamma$ is the shaping parameter and $r$ is the support ratio.

Direction assignment:
\[
\text{if } o < 0: (I_L, I_R) = (I_d, I_s), \qquad
\text{else}: (I_L, I_R) = (I_s, I_d)
\]

Each motor is updated continuously in real-time playback mode using an 8-bit command value:
\[
v = \lfloor \text{clip}(I, 0, 1) \cdot 127 \rfloor
\]

This produces smooth vibration transitions as the target moves across the frame, avoiding abrupt switching and improving directional guidance clarity.

\subsection{Next Steps and Iterative Refinement}

Empirical feedback gathered directly from qualitative responses asserts definitively that the continuous gradient haptic mapping algorithms exhibit vastly superior interpretative thresholds compared to original distinct pulse mapping modes. In moving toward finalization, the subsequent iteration phases aim to consolidate hardware footprint for superior wearable ergonomics, refine model quantization to aggressively limit structural end-to-end latency, and expand vocabulary mapping accuracy. The ultimate goal is delivering a comprehensively safe, predictable, and fully embedded multi-modal navigation system.
 
\clearpage
\appendix
\section{Appendix}

\subsection{Appendix A: Template Inclusions}

% NOTE: The following bubble sort example is retained strictly to fulfill the original template and assignment guidelines requiring theoretical computational algorithmic representation.

\begin{algorithm}[h]
   \caption{Bubble Sort}
   \label{alg:example}
\begin{algorithmic}
   \STATE {\bfseries Input:} data $x_i$, size $m$
   \REPEAT
   \STATE Initialize $noChange = true$.
   \FOR{$i=1$ {\bfseries to} $m-1$}
   \IF{$x_i > x_{i+1}$}
   \STATE Swap $x_i$ and $x_{i+1}$
   \STATE $noChange = false$
   \ENDIF
   \ENDFOR
   \UNTIL{$noChange$ is $true$}
\end{algorithmic}
\end{algorithm}

\subsection{Appendix B: Project Timeline and Task Management}

Figure~\ref{fig:timeline} illustrates the master timeline and milestone tracking documentation maintained via Google Docs. This shared tracking framework ensured synchronous cross-team development and effectively aligned both hardware and software delivery schedules. 

\begin{figure*}[h]
    \centering
    \includegraphics[width=\textwidth]{figs/timeline.png}
    \caption{Project timeline and task distribution maintained via Google Docs.}
    \label{fig:timeline}
\end{figure*}

\subsection{Appendix C: System Visualizer Interface}

Figure~\ref{fig:website} displays the lightweight web-based observability interface built to debug the live inference loop. The visualizer streams the on-device camera feed, YOLO-World detection bounding boxes, and internal system state transitions to assist in rapidly diagnosing latency and localization deviations.

\begin{figure*}[h]
    \centering
    \includegraphics[width=\textwidth]{figs/website.png}
    \caption{The diagnostic observability interface capturing live camera feeds and detection confidence.}
    \label{fig:website}
\end{figure*}

\clearpage
\onecolumn
\bibliography{example-ref}

\end{document} 

